\newpage
\section{Objetivos específicos}
\label{sec:objetivos}

O objetivo geral deste trabalho é desenvolver um aplicativo móvel de
treinamento individual de forró universitário que utiliza visão computacional
para analisar em tempo real a transferência de peso, a postura corporal, o
ritmo e a execução de passos do praticante, fornecendo feedback imediato sem
a necessidade de um instrutor presencial. Para alcançar esse objetivo, foram
definidos os seguintes objetivos específicos:

\begin{enumerate}
    \item Sistematizar, a partir da revisão bibliográfica, os critérios
    biomecânicos relevantes para avaliação dos quatro módulos do sistema ---
    transferência de peso, postura, ritmo e execução de passos do forró
    universitário --- definindo as métricas e os limiares de avaliação a
    serem adotados;

    \item Projetar a arquitetura do sistema Dança~AI, especificando os quatro
    módulos de análise, o pipeline de captura e inferência (CameraX →
    MediaPipe Pose Landmarker → módulos → feedback) e as interfaces entre
    componentes;

    \item Implementar o módulo de transferência de peso, classificando o
    apoio do praticante como esquerdo, direito ou neutro a partir da razão
    horizontal entre os landmarks de quadril detectados pelo MediaPipe, com
    feedback visual simples e direto;

    \item Implementar o módulo de análise de postura, avaliando o
    alinhamento do tronco e dos ombros por meio de regras geométricas sobre
    os landmarks da parte superior do corpo e de um modelo de classificação
    embarcado treinado sobre coordenadas $xyz$ de dançarinos, com feedback visual indicando
    conformidade ou desvio postural;

    \item Implementar o módulo de ritmo, com metrônomo configurável por BPM e
    tipo de batida, feedback auditivo ao praticante e validação rítmica por
    comparação entre a cadência do metrônomo e o padrão de transferência de
    peso detectado em tempo real;

    \item Implementar o módulo de identificação de movimentos, com
    classificador embarcado treinado para reconhecer os seis passos definidos
    (Balanço, Base~1, Base~2, Abertura, Giro Invertido e Giro em Linha), nos
    modos Livre --- o sistema identifica e avalia o passo executado --- e
    Desafio --- o praticante deve seguir uma sequência predefinida;

    \item Avaliar o sistema desenvolvido em termos de latência do pipeline
    completo, acurácia dos modelos de classificação e experiência do usuário,
    por meio de testes com praticantes de forró utilizando a escala SUS
    (\textit{System Usability Scale}).
\end{enumerate}

A execução desses objetivos está distribuída entre as duas etapas do
trabalho de conclusão de curso: a presente etapa (TCC01) compreende a
fundamentação teórica e o projeto da arquitetura (objetivos 1 e 2), enquanto
o desenvolvimento, integração e avaliação dos módulos (objetivos 3 a 7) serão
realizados na etapa subsequente (TCC02).
